\documentclass{IEEEtran}
\interdisplaylinepenalty=2500

\usepackage{tikz}
\usetikzlibrary{arrows.meta}
% ----- Encoding & language -----
\usepackage[utf8]{inputenc}
\usepackage[T5]{fontenc}
\usepackage[vietnamese]{babel}

% ----- Math & typography -----
\usepackage{amsmath,amssymb,amsfonts}
\usepackage{microtype}            % cải thiện giãn chữ, giảm overfull
\allowdisplaybreaks               % cho phép tách cụm phương trình qua cột/trang

% (Tuỳ chọn) Ký hiệu nhanh cho tập số thực
\newcommand{\R}{\mathbb{R}}
\microtypesetup{protrusion=false}
\usepackage{float}
\usepackage{tikz}
\usetikzlibrary{arrows.meta,positioning,calc}
\usepackage{booktabs} % preamble
\usepackage[caption=false,font=footnotesize]{subfig}


\usepackage{siunitx}
\sisetup{
  table-number-alignment = center,
  round-mode = places,
  round-precision = 2,
  scientific-notation = true
}

% ===== Title =====
\title{SphyL-Vessel: Ràng buộc vật lý lồi--trơn cho dự báo quỹ đạo tàu từ dữ liệu AIS}

\author{
  Nhữ Phạm Quang Mạnh\\
  Nghiên cứu độc lập\\
  TP. Hồ Chí Minh, Việt Nam\\
  \texttt{py.quangmanh.ai@gmail.com}
}

\begin{document}
\maketitle

\begin{abstract}
Dự báo quỹ đạo tàu biển là bài toán trọng yếu cho an toàn hàng hải, quản lý luồng giao thông và điều hướng tự động. 
Các mô hình học sâu gần đây như LSTM, Transformer hay diffusion đã cải thiện độ chính xác, 
nhưng thường đòi hỏi chi phí tính toán lớn, thiếu ràng buộc vật lý để đảm bảo khả dĩ, 
và khó diễn giải trong ứng dụng. 

Chúng tôi giới thiệu \emph{SphyL--Vessel}: một khung \emph{loss vật lý lồi--trơn} có thể tích hợp trực tiếp vào bất kỳ mô hình tuần tự hiện có, 
không cần thay đổi kiến trúc hay tăng chi phí tính toán. 
Ý tưởng cốt lõi là mô hình hóa miền khả đạt trong một bước thời gian dưới dạng hình học lồi 
và áp dụng phạt mềm bằng \textit{softplus}, giúp duy trì gradient mượt và tối ưu ổn định. 
Từ elip cơ bản, chúng tôi mở rộng thành siêu elip dịch chuyển cho tình huống rẽ/phanh gắt 
và hành lang bất đối xứng (\emph{tube--stomach}) cho chuyển động phức tạp, 
tất cả đều bảo toàn lồi tính toàn cục. 

Thực nghiệm trên dữ liệu AIS California (04/2023) cho thấy SphyL--Vessel duy trì sai số RMSE ngang với baseline, 
nhưng nâng cao đáng kể độ phủ (coverage) và tỉ lệ trúng trong bán kính an toàn (Hit@100m). 
Kết quả này khẳng định rằng một tiếp cận nhẹ nhưng gắn chặt vật lý có thể vừa khả dĩ, 
vừa hiệu quả hơn các mô hình dự báo tàu biển hiện nay.
\end{abstract}


\begin{IEEEkeywords}
AIS, dự báo quỹ đạo, ràng buộc vật lý, tối ưu lồi, super-ellipse, corridor, softplus hinge
\end{IEEEkeywords}


\section{Nghiên cứu liên quan}
Bài toán dự báo quỹ đạo tàu biển có thể chia thành bốn nhóm chính: 
(i) mô hình thống kê truyền thống, 
(ii) mô hình học sâu tuần tự, 
(iii) kiến trúc tiên tiến như Transformer, đồ thị và mô hình sinh, 
và (iv) phương pháp tích hợp ràng buộc vật lý hoặc tri thức miền.

\subsection{Mô hình thống kê truyền thống}
Các phương pháp kinh điển như Kalman filter~\cite{kalman1960} và ARIMA~\cite{box1970} là nền tảng ban đầu cho dự báo chuỗi thời gian, 
có ưu điểm dễ diễn giải và nhẹ tính toán. 
Trong dữ liệu AIS, Kalman filter thường được áp dụng để lọc nhiễu tọa độ GPS và nội suy quỹ đạo, 
trong khi ARIMA mô hình hóa xu hướng tốc độ hoặc hướng đi theo chuỗi một chiều.
Tuy nhiên, cả hai đều giả định tuyến tính và dừng yếu, khiến chúng không phù hợp 
với hành vi phi tuyến của tàu trong kịch bản cua gắt, phanh hoặc thay đổi luồng và tương tác phức tạp trong dữ liệu AIS~\cite{chen2023ais,li2023ml}.

\subsection{Mô hình học sâu tuần tự}
RNN và LSTM tận dụng tốt tính tuần tự của dữ liệu AIS (vị trí, vận tốc, hướng ghi nhận mỗi 60 giây), 
cho kết quả chính xác hơn hẳn so với mô hình tuyến tính trong điều kiện hành trình ổn định. 
Nhiều nghiên cứu đã bổ sung cơ chế attention hoặc kiến trúc encoder--decoder 
để mở rộng, cải thiện chân trời dự báo lên nhiều bước~\cite{li2023deepais,zhang2022lstm,li2023engapp}. 
Tuy vậy, do tối ưu hoàn toàn trên dữ liệu mà không có ràng buộc vật lý, 
các mô hình này có thể sinh ra quỹ đạo phi thực tế---tăng tốc đột ngột, 
hoặc thay đổi hướng vượt giới hạn động học của tàu.

\subsection{Transformer, đồ thị và mô hình sinh}
Transformer cho AIS (ví dụ TrAISformer) thể hiện khả năng mô hình hóa phụ thuộc dài hạn và đa mô thức~\cite{wang2024transformer,wang2024isprs}. 
Các mô hình đồ thị không gian--thời gian và diffusion cho phép nắm bắt tương tác đa tác nhân phức tạp~\cite{liu2023graphais,arxiv2023multipath}. 
Tuy nhiên, những hướng này thường đòi hỏi chi phí tính toán lớn, khó triển khai trong thời gian thực.

\subsection{Ràng buộc vật lý và tri thức miền}
Các tiếp cận physics-informed như PINN đưa phương trình động học vào loss qua các sơ đồ rời rạc (Euler, midpoint)~\cite{raissi2019pinn}. 
TPK (Trustworthy Prior Knowledge) khai thác tri thức miền và giới hạn động học để bảo đảm quỹ đạo khả thi hơn~\cite{zhang2023tpk}. 
Ngoài ra, một số nghiên cứu gần đây kết hợp deep learning với tri thức động lực học hoặc framework đa mô thức~\cite{zaman2024plos,li2024maker}. 
Khác với các hướng trên, \emph{SphyL--Vessel} đi theo con đường hình học lồi: 
ràng buộc được mô hình hóa bằng miền elip hoặc hành lang "tube--stomach", 
giữ loss đơn giản, dễ tích hợp và thân thiện tính toán, đồng thời mở rộng được sang nhiều dạng hình học khác.

\textbf{Lượng physics loss của hướng tiếp cận SphyL-Vessel đơn giản hơn các hướng tiếp cận physics trước:}
Các phương pháp trước (PINN, TPK) và SphyL--Vessel đều đưa tri thức vật lý vào hàm mất mát,
nhưng có sự khác biệt căn bản ở ba khía cạnh:

\emph{(a) Dạng toán học của ràng buộc.}
PINN yêu cầu rời rạc hóa phương trình vi phân (ví dụ sơ đồ Euler hoặc midpoint):
\[
  \mathcal{L}_{\text{PINN}} = \bigl\|\dot{x} - f(x,t;\theta)\bigr\|^2,
\]
trong đó $f$ là vế phải của phương trình động học, đòi hỏi tính đạo hàm 
bậc hai theo cả tham số mô hình lẫn thời gian~\cite{raissi2019pinn}.
Tương tự, TPK xây dựng bất đẳng thức kiểm tra cho từng ràng buộc động học riêng lẻ 
(tốc độ góc, gia tốc dọc), dẫn đến một hệ nhiều điều kiện phải thỏa đồng thời~\cite{zhang2023tpk}.
Ngược lại, SphyL--Vessel gộp \emph{toàn bộ} không gian khả đạt vào một bất đẳng thức duy nhất:
\[
  \phi_{\text{lat/fwd}} = \bigl[\,g(s,y)\,\bigr]_+^2,
\]
trong đó $g(s,y)$ là hàm hình học lồi--trơn (quadratic form elip hoặc biên hành lang).
Mỗi bước dự báo chỉ cần một phép nhân vô hướng và một lần gọi \textit{softplus},
chi phí tính toán là $\mathcal{O}(1)$ so với kích thước mô hình.

\emph{(b) Chiều sâu tích hợp vào kiến trúc.}
PINN thường yêu cầu mạng học xấp xỉ cả nghiệm lẫn đạo hàm của phương trình,
nên buộc phải thiết kế lại kiến trúc hoặc dùng automatic differentiation liên tục
qua nhiều bước thời gian---không tương thích trực tiếp với LSTM hay Transformer.
TPK cần gắn kèm module kiểm tra sau mỗi bước dự báo.
Trong khi đó, $\mathcal{L}_{\text{phys}}$ của SphyL--Vessel
là \emph{kiến trúc--bất biến}: bất kỳ mô hình nào có đầu ra dạng $(s, y)$
đều có thể cộng thêm loss này mà không cần sửa đổi nội bộ.

\emph{(c) Độ mượt gradient và ổn định huấn luyện.}
Các bất đẳng thức cứng (hard constraint) trong TPK hoặc điều kiện \textit{residual} 
trong PINN dễ tạo ra gradient suy biến (vanishing/exploding) ở biên ràng buộc.
SphyL--Vessel thay hinge cứng bằng \textit{softplus--hinge}
$\phi_\kappa(u) = \tfrac{1}{\kappa}\log(1+e^{\kappa u})$,
bảo đảm gradient liên tục và khác không ở mọi điểm, giúp tối ưu hội tụ mượt mà hơn.
Ngoài ra, lịch ramp tuyến tính $\lambda(t)$ tăng dần từ 0 ngăn mất mát vật lý
áp đảo mất mát dữ liệu ở giai đoạn đầu huấn luyện---một kỹ thuật không cần thiết
nếu loss vật lý đã đủ đơn giản và trơn từ đầu.

Tóm lại, ưu điểm cốt lõi của SphyL là đại số hóa ràng buộc hình học:
thay vì giải phương trình vi phân hay kiểm tra từng bất đẳng thức động học,
toàn bộ vùng khả đạt được biểu diễn bằng một hàm lồi duy nhất,
vừa đủ đơn giản để tích hợp, vừa đủ giàu để phản ánh vật lý thực tế.

\begin{table}[H]
\caption{So sánh tổng quan các hướng tiếp cận dự báo quỹ đạo tàu.}
\label{tab:related-compare}
\centering
\renewcommand{\arraystretch}{1.15}
\begin{tabular}{lcccc}
\hline
\textbf{Phương pháp} & \textbf{Độ chính xác} & \textbf{Ràng buộc} & \textbf{Chi phí} & \textbf{Tích hợp} \\
                    &                       & \textbf{vật lý}   & \textbf{tính toán} & \textbf{kiến trúc} \\
\hline
Kalman / ARIMA      & Thấp                  & Không              & Rất nhẹ          & --           \\
LSTM / RNN         & Trung bình            & Không              & Nhẹ              & Độc lập      \\
Transformer/Diffusion & Cao                 & Không              & Nặng             & Phụ thuộc   \\
PINN               & Cao                   & Có (PDE)           & Nặng             & Cần thiết kế lại \\
TPK                & Trung--Cao            & Có (bất đẳng thức) & Trung bình       & Module riêng \\
\textbf{SphyL--Vessel} & \textbf{Trung--Cao} & \textbf{Có (hình học lồi)} & \textbf{Rất nhẹ} & \textbf{Bất biến} \\
\hline
\end{tabular}
\end{table}


\subsection{ Bộ ký hiệu và đại lượng vật lý}
\begin{table}[H]
\caption{Các ký hiệu chính trong \emph{SphyL--Vessel}.}
\label{tab:sphyl-symbols}
\centering
\begin{tabular}{ll}
\hline
Ký hiệu & Diễn giải \\
\hline\\[4pt]
$p_0=(x_0,y_0)$ & Vị trí ban đầu trong hệ toàn cục \\
$\hat{p}=(x,y)$ & Vị trí dự đoán sau khoảng $\Delta t$ \\
$v_0$ & Vận tốc tại $t_0$ \\
$\theta_0$ & Góc heading ban đầu so với trục $Ox$ \\
$d=\theta_{\text{motion}}-\theta_0$ & Chênh lệch hướng chuyển động và heading \\
$\Delta t$ & Độ dài bước dự báo \\
$a_t^{\max}$ & Gia tốc cực đại dọc (thúc--phanh) \\
$a_n^{\max}$ & Gia tốc cực đại ngang (rẽ) \\
$v_{\max}$ & Vận tốc cực đại \\
$e_{\parallel},e_{\perp}$ & Cơ sở trực chuẩn cục bộ (song song/vuông góc) \\
$\Delta p=\hat{p}-p_0$ & Vectơ dịch chuyển \\
$s=\Delta p \cdot e_{\parallel}$ & Tiến dọc \\
$y=\Delta p \cdot e_{\perp}$ & Lệch ngang \\
$a,b$ & Bán trục dọc/ngang của elip \\
$M$ & Ma trận xác định dương trong quadratic form \\
$\tau$ & Tham số kích hoạt mềm (turn-rate/gia tốc) \\
$r_{\text{base}}, B_{\max}, s_{\max}$ & Tham số hành lang tube \\
$\rho(s)$ & Hàm bụng chuẩn $4t(1-t)$ \\
$\rho_{\alpha,\beta}(t)$ & Hàm bụng cải tiến Beta-shape \\
$p_c(s)$ & Đường centerline cong \\
$s_{\text{cap}}$ & Giới hạn động (forward cap) \\
$\phi_{\text{lat}},\phi_{\text{fwd}},\phi_{\text{prog}}$ & Phạt loss ngang/dọc/tiến \\
$\lambda, \alpha_m$ & Hệ số loss vật lý / pha trộn dữ liệu \\
$\phi_{\kappa}(u)$ & Softplus: $\tfrac{1}{\kappa}\log(1+e^{\kappa u})$ \\
\\[4pt]
\hline
\end{tabular}
\end{table}


\subsection{ Khung toạ độ cục bộ}
Từ trạng thái ban đầu, thiết lập hệ trục:
\[
e_{\parallel} = R(d)\,e_{\text{head}}, 
\qquad 
e_{\perp} = (-e_{\parallel,y},\,e_{\parallel,x}),
\]
trong đó $R(d)$ là phép quay một góc $d$.  
Mọi dự đoán $\hat{p}$ được quy chiếu về tọa độ cục bộ:
\[
s=(\hat{p}-p_0)\cdot e_{\parallel}, 
\qquad 
y=(\hat{p}-p_0)\cdot e_{\perp}.
\]

\textbf{Centerline cong.}  
Trong phiên bản cải tiến, thay vì giữ trục $e_{\parallel}$ cố định, 
chúng tôi định nghĩa đường trung tâm cong:
\[
p_c(s) = p_0 + s e_{\parallel} + \delta(s)\,e_{\perp}, 
\qquad \delta(s) = \alpha_{\text{curve}}\,s(1-s),
\]
cho phép khung toạ độ $(s,y)$ uốn theo hướng chuyển động, 
phản ánh động học tốt hơn khi tàu cua hoặc đổi hướng.

\subsection{ Miền elip cứng: khởi nguyên lý thuyết}
Trong bước thời gian $\Delta t$, quỹ đạo khả đạt giới hạn bởi:
\begin{IEEEeqnarray}{rCl}
0 \le s &\le& v_0 \Delta t + \tfrac12 a_t^{\max}\,\Delta t^{2},\\
|y|     &\le& \tfrac12 a_n^{\max}\,\Delta t^{2},\\
v_0     &\le& v_{\max}.
\end{IEEEeqnarray}
Hình học hóa: mô hình miền khả đạt là elip centered tại $p_0$, với bán trục
\[
a = v_0 \Delta t + \tfrac{1}{2} a_t^{\max} (\Delta t)^2, 
\qquad 
b = \tfrac{1}{2} a_n^{\max} (\Delta t)^2.
\]
Điều kiện khả thi:
\[
\mathcal{E} = \{(s,y) \;\mid\; \tfrac{s^2}{a^2} + \tfrac{y^2}{b^2} \leq 1 \}.
\]

\subsection{ Elip dịch tâm: bù lệch hình học}
Để khắc phục sự đối xứng cứng nhắc, ta cho phép tâm elip dịch chuyển theo hướng tiến dọc. 
Trực giác: khi tàu cua hoặc phanh gắt, vùng khả thi thực tế thường dồn về phía trước (do quán tính) hoặc lệch sang trong cua.  
Miền mới:
\[
\mathcal{E}_{\text{shift}} = \Big\{ (s,y) \;\Big|\; 
\tfrac{(s-c_s)^2}{a^2} + \tfrac{(y-c_y)^2}{b^2} \leq 1 \Big\}.
\]

\subsection{ Siêu elip: mở rộng linh hoạt}
Một mở rộng tự nhiên là thay norm $L_2$ bằng $L_p$ với $p>2$, 
giúp biên cong phẳng hơn và nội dung bao quát rộng hơn:
\[
\mathcal{E}_{p} = \Big\{ (s,y) \;\Big|\; 
\Big(\tfrac{|s|}{a}\Big)^p + \Big(\tfrac{|y|}{b}\Big)^p \leq 1 \Big\}.
\]
Siêu elip vẫn giữ tính lồi (với $p \geq 2$) nhưng giảm hiện tượng ``dẹt'' khi $b$ nhỏ.

\subsection{ Tube--Stomach Corridor: hành lang bất đối xứng}
Khắc phục các hạn chế, hành lang được định nghĩa theo $s$:
\[
0 \leq s \leq s_{\max}, 
\qquad 
-|y_R(s)| \leq y \leq y_L(s).
\]

\textbf{Phiên bản gốc.}  
Với $\rho(s)=4t(1-t),\,t=s/s_{\max}$, 
biên ngang:
\[
y_L(s) = r_{\text{base}} + \tau \cdot B(s), \qquad 
y_R(s) = r_{\text{base}} + \tau \cdot B(s).
\]

\textbf{Phiên bản cải tiến.}  
Thay vì dùng $\rho(s)$ cứng nhắc, 
chúng tôi áp dụng hàm bụng Beta-shape:
\[
\rho_{\alpha,\beta}(t) = 
\frac{t^{\alpha}(1-t)^{\beta}}
     {\left(\tfrac{\alpha}{\alpha+\beta}\right)^{\alpha}
      \left(\tfrac{\beta}{\alpha+\beta}\right)^{\beta}}, 
\qquad t = s/s_{\max}.
\]
Hàm này cho phép điều chỉnh linh hoạt độ phình, 
bao quát tốt hơn khi cua gắt hoặc phanh mạnh.

\subsection{ Loss lồi--trơn và cơ chế tối ưu}
Loss tổng quát gồm ba thành phần:
\[
\mathcal{L} = (1-\alpha_m)\,\mathcal{L}_{\text{MSE}}
+ \alpha_m \,\mathcal{L}_{\text{MSE-meters}}
+ \lambda \,\mathcal{L}_{\text{phys}}.
\]

\textbf{Cải tiến cap động \& progress anchor.}  
Để ngăn dự đoán phi lý:
\[
s_{\text{cap}} = \text{clip}\Big((1+0.3\tau)\,s_{\text{tgt}},\,0.3\,s_{\max},\,1.6\,s_{\max}\Big),
\]
và bổ sung phạt tiến đủ:
\[
\phi_{\text{prog}} = [s_{\text{tgt}}-s_{\text{pred}}]_+^2.
\]

\textbf{Hình thức cuối cùng.}  
\[
\mathcal{L}_{\text{phys}} = 
w_{\text{lat}}\phi_{\text{lat}} +
w_{\text{fwd}}\phi_{\text{fwd}} +
w_{\text{prog}}\phi_{\text{prog}},
\]
trong đó
\[
\phi_{\text{lat}} = [|y|-y_{\text{bdry}}(s)]_+^2,\quad
\phi_{\text{fwd}} = [s_{\text{pred}}-s_{\text{cap}}]_+^2.
\]
Mọi phạt đều qua softplus $\phi_\kappa(u)$ để duy trì gradient mượt.


\subsection{ Thảo luận các biến thể hình học}

\begin{figure}[H]
  \centering
  % ảnh elip (trên) — phóng to thêm 10%
  \scalebox{1.0}{\includegraphics[width=0.9\columnwidth]{tube.pdf}}\\[1ex]

  % ảnh tube-stomach (dưới)
  \includegraphics[width=0.9\columnwidth]{tube_3.pdf}
  \caption{So sánh hai dạng vùng khả dĩ: (trên) Feasible ellipse và (dưới) Tube--stomach.}
  \label{fig:ellipse-stomach}
\end{figure}



\textbf{1) Elip cứng.}  
Mô hình elip đơn giản, quadratic form gọn gàng, chứng minh tính lồi tức thì. 
Đây là \emph{khởi nguyên lý thuyết} của SphyL--Vessel. 
Tuy nhiên, với dữ liệu AIS, $b \approx 0$ khiến elip bị dẹt; 
chỉ cần nhiễu nhỏ đã rơi ra ngoài, dẫn tới loss cực lớn. 
Do đó, elip cứng đẹp về mặt hình học nhưng mong manh trong thực nghiệm.  

\textbf{2) Elip dịch tâm.}  
Cho phép dịch tâm theo hướng tiến hoặc vào cua để khắc phục phần nào tính đối xứng. 
Biến thể này mô tả tốt hơn khi tàu quay hoặc phanh, 
nhưng vẫn chỉ là affine transform của elip, 
chưa có ``bụng'' để bao quát quán tính trong các tình huống cực trị.  

\textbf{3) Siêu elip.}  
Với $p>2$, biên cong phẳng hơn, coverage và Hit@100m cải thiện rõ rệt. 
Siêu elip vẫn là tập lồi (convex set) và dễ tối ưu, 
nhưng chưa có cơ chế động để thích nghi với trạng thái cua hoặc giảm tốc.  

\textbf{4) Tube--Stomach.}  
Bước nhảy vọt hình học: bất đối xứng tự nhiên, $s_{\max}$ động theo $v_0$, 
biên ngang $r_{\text{base}}+B_{\max}\rho(s)$ tạo ``bụng'' mềm để bao quát quán tính khi cua/phanh. 
Trong phiên bản cải tiến, hàm $\rho(s)$ được thay bằng $\rho_{\alpha,\beta}(t)$ dạng Beta-shape 
và kết hợp với centerline cong, giúp thích nghi tốt hơn với quỹ đạo thực tế.  
Tuy nhiên, khác với elip/siêu elip, hình học tube--stomach không còn là tập lồi toàn cục: 
đoạn thẳng nối hai điểm trong hành lang có thể chọc ra ngoài ở giữa. 
Dù vậy, nhờ thiết kế phạt mềm bằng hinge và softplus, 
loss vẫn giữ được đặc tính tối ưu mượt và ổn định, 
trở thành lựa chọn thực tế hơn trong huấn luyện.  

\medskip
\noindent \textbf{Tóm lại:} 
elip cứng đặt nền tảng lý thuyết, 
elip dịch tâm và siêu elip mở rộng tính bao phủ, 
còn tube--stomach (với bụng linh hoạt và centerline cong) 
mang lại giải pháp gần gũi thực tế nhất. 
Các biến thể này phản ánh lộ trình tiến hoá từ đơn giản $\to$ thích nghi $\to$ giàu động học.



% hinh tik
% hinh tik

\section{Thảo luận}

\subsection{So sánh tổng quan các biến thể}
Các thí nghiệm cho thấy baseline có thể chấp nhận được ở offshore (RMSE $\sim$1029 m, Hit@100 $\sim$0.42), 
nhưng sụp đổ trong vịnh San Francisco (Hit@100 chỉ 0.28). 
Ngược lại, khi tích hợp ràng buộc hình học, 
Ellipse-only $\tau$-adapt và Tube--Stomach đều phục hồi Hit@100 lên 0.565, 
giảm mạnh RMSE/MAE/Median. 
Điều này chứng minh rằng lợi ích của ràng buộc vật lý càng rõ rệt trong môi trường có động lực phức tạp.

\subsection{Elip: chặt chẽ và bao phủ}
Hình học elip duy trì coverage cao, phù hợp với kịch bản tàu đi thẳng tốc độ cao. 
Nhờ $\tau$-adapt, bán trục $a, b$ được điều chỉnh linh hoạt theo turn-rate và gia tốc thực tế, 
nên ellipse cũng thích nghi tốt khi có cua nhẹ.
Tuy nhiên, trong môi trường ngoài khơi với nhiễu GPS lớn, biên elip hẹp khiến 
các dự báo lệch nhỏ cũng bị phạt nặng---dẫn đến outside severity rất lớn. 
Đây là đánh đổi cố hữu của hình học chặt: coverage cao nhưng kém khoan dung trước nhiễu.

\subsection{Tube--Stomach: thực dụng và ổn định}
Tube--Stomach với bụng Beta-shape, centerline cong và cap động 
tạo hành lang rộng rãi, ổn định hơn khi tàu cua hoặc phanh gấp. 
Cụ thể, hàm $\rho_{\alpha,\beta}(t)$ làm phồng biên ngang ở đoạn giữa hành lang---
nơi quán tính tàu đạt đỉnh---nên bao phủ được các quỹ đạo lệch do cua gắt mà elip cứng bỏ sót.
Fairness analysis chỉ ra outside severity $\approx$0, median và inside margin P90 ổn định, 
xác nhận hành lang bao phủ sát thực tế mà không quá chặt. 
Coverage tuyệt đối thấp hơn ellipse do hành lang rộng hơn, 
nhưng đổi lại mô hình ít bị phạt oan, ổn định hơn trong huấn luyện và khả năng triển khai thực tế cao hơn.

\subsection{Bài học thực nghiệm}
Bức tranh chung:
\begin{itemize}
  \item Ellipse thắng về coverage, đặc biệt ở offshore.
  \item Tube thắng về fairness và ổn định, đặc biệt trong vịnh San Francisco.
\end{itemize}
Do đó, việc chọn biến thể phụ thuộc bối cảnh:
safety-critical cần coverage $\to$ ellipse;
ứng dụng thực tiễn trong môi trường phức tạp $\to$ tube.
Điểm chung của cả hai là chi phí tính toán cận như không đổi so với baseline thuần,
nhưng mang lại cải thiện đáng kể về tính khả dĩ vật lý.

\subsection{Hạn chế và hướng mở rộng}
Mặc dù đạt được kết quả khả quan, SphyL--Vessel còn một số hạn chế cần thừa nhận.
Thứ nhất, ngưỡng vật lý ($a_t^{\max}$, $a_n^{\max}$, $v_{\max}$) được ước lượng toàn cục 
theo loại tàu, chưa thích nghi theo thời gian thực với điều kiện thời tiết hay tải trọng.
Thứ hai, công trình chỉ đánh giá dự báo một bước; 
khi mở rộng sang đa bước (multi-step), sai số tích lũy có thể làm suy giảm hiệu quả ràng buộc.
Thứ ba, tương tác đa tàu---quan trọng trong vùng cảng đông đúc---chưa được mô hình hóa;
cần kết hợp thêm ràng buộc tránh va chạm hoặc biểu diễn đồ thị. 
Các hướng này sẽ được khám phá trong nghiên cứu tương lai.





\bibliographystyle{IEEEtran}
\bibliography{refs}
\end{document}